\section{Methodology}\label{sec:methodology}

\subsection{Template Design Philosophy}

The Seka template follows three core principles:

\begin{enumerate}
    \item \textbf{Simplicity:} No special fonts or external dependencies
    \item \textbf{Modularity:} Sections stored in separate files for organization
    \item \textbf{Professionalism:} Academic formatting standards maintained
\end{enumerate}

\subsection{Document Class Configuration}

The base class is \texttt{extarticle} with options for 10pt font, two-column layout, and two-sided printing. The geometry package configures margins for optimal readability:

\begin{equation}\label{eq:margins}
\begin{aligned}
\text{top margin} &= 2.2 \text{ cm} \\
\text{bottom margin} &= 2.2 \text{ cm} \\
\text{left margin} &= 1.8 \text{ cm} \\
\text{right margin} &= 1.8 \text{ cm}
\end{aligned}
\end{equation}

\subsection{Color Scheme}

An academic color palette enhances visual hierarchy without distraction:

\begin{itemize}
    \item \textcolor{titlecolor}{Title color: Navy blue} for main title
    \item \textcolor{sectioncolor}{Section color: Dark blue} for headings
    \item \textcolor{linkcolor}{Link color: Medium blue} for hyperlinks
\end{itemize}

\subsection{Custom Environments}

Three custom environments improve document structure:

\begin{enumerate}
    \item \texttt{note}: For supplementary comments (light blue background)
    \item \texttt{important}: For critical information (orange highlight)
    \item \texttt{code}: For code listings with proper formatting
\end{enumerate}

Equation~\ref{eq:custom} demonstrates mathematical formatting.

\begin{equation}\label{eq:custom}
\mathcal{L}(\theta) = \prod_{i=1}^{n} p(x_i | \theta)
\end{equation}

\begin{important}
Always cite sources properly when referencing external work. Use \verb|\cite{}| for individual citations or \verb|\citep{}| for parenthetical citations.
\end{important}